\chapter{Source Code and Reproducibility}
\label{app:reproducibility}

The complete source code of the \textit{Omnideps} framework,
including the heuristic extraction engine, the resolution algebra,
the embedded web visualizer, and the experimental benchmark suite
developed for this thesis, is open-source and publicly available.

\section{Repository Availability}
The project repository can be accessed at the following URL:
\begin{center}
  \url{https://github.com/FerrarioChristian/omnideps}
\end{center}

\section{Installation for End Users}
\textit{Omnideps} is distributed as a single standalone binary. The
web visualizer (SvelteKit) is pre-compiled and embedded directly
inside the Rust executable. To use the tool, it is not necessary to
install Node.js, npm, or Rust.
Users can simply:
\begin{enumerate}
  \item Navigate to the \textbf{Releases} page on the GitHub repository.
  \item Download the executable for their Operating System (Windows
    \texttt{.exe}, macOS, or Linux).
  \item Run it directly from the terminal.
\end{enumerate}

\section{Command Line Interface}
The framework provides several subcommands to orchestrate the
analysis and visualization process.

\subsection*{Analyze Code}
To extract the dependency graph from a target project, use the
\texttt{analyze} command:
\begin{verbatim}
omnideps analyze /path/to/project --output graph.json
\end{verbatim}
This command will parse the codebase, execute the extraction
heuristics, run the resolution engine, and output the final
dependency graph in standard JSON format.

\subsection*{Custom Configuration}
To initialize a default configuration file (\texttt{omnideps.json})
that can be customized to override the default extraction rules and
resolution strategies, run:
\begin{verbatim}
omnideps config init
\end{verbatim}
The generated configuration file can then be provided to the analyzer
using the \texttt{--config} flag (e.g., \texttt{omnideps analyze
/path --config omnideps.json}).

\subsection*{Web Visualizer}
To interactively explore the generated graphs, the embedded web
interface can be launched via:
\begin{verbatim}
omnideps serve
\end{verbatim}
This starts an ultra-lightweight local HTTP server (defaulting to
port 3000) and automatically opens the default web browser to the
interactive UI.

\subsection*{Benchmarking Suite}
The project includes custom benchmarks to accurately measure false
positives and false negatives on the AST abstraction across the
supported languages.
To run a specific benchmark:
\begin{verbatim}
omnideps benchmark run tests/benchmarks/benchmark-rust -o /output_dir
\end{verbatim}
To execute the entire suite and generate the aggregated CSV report:
\begin{verbatim}
omnideps benchmark all -o /output_dir
\end{verbatim}

\section{Building from Source}
For development and reproducibility purposes, the project can be
built from source. This requires both \textbf{Node.js} and the
\textbf{Rust Toolchain} (\texttt{cargo} and \texttt{rustc}).

First, build the SvelteKit frontend, which generates static files in
the \texttt{visualizer-svelte/build} directory:
\begin{verbatim}
cd visualizer-svelte
npm install
npm run build
cd ..
\end{verbatim}

Next, compile the Rust executable. The \texttt{rust-embed} macro will
automatically package the compiled Svelte web assets inside the
binary during compilation:
\begin{verbatim}
cargo build --release
\end{verbatim}
The final executable will be located in \texttt{target/release/omnideps}.

\section{Testing}
The framework includes a comprehensive suite of unit and integration
tests designed to validate the internal extraction and resolution
logic. To execute the test suite, run:
\begin{verbatim}
cargo test
\end{verbatim}
